% =============================================================================
\section{Dimension 3: Pool Architecture}
\label{sec:pool}
% =============================================================================

Our fleet optimization tools~\cite{fleetopt2026, fleetsim2026,
onewlaw2026} provide the analytical and simulation models for pool
architecture.

% -----------------------------------------------------------------------------
\subsection{Homogeneous vs.\ Heterogeneous GPU Pools}
\label{sec:pool:homo-hetero}
% -----------------------------------------------------------------------------

FleetOpt~\cite{fleetopt2026} derives the minimum-cost fleet
analytically: a two-pool architecture with optimal boundary
$\boundary^*$ satisfying an equal marginal GPU cost condition.  On
production traces, this reduces GPU cost by 6--82\% versus
homogeneous fleets.  inference-fleet-sim~\cite{fleetsim2026} extends
this to heterogeneous GPU types (A10G, A100, H100) with a
physics-informed performance model.

M\'{e}lange~\cite{griggs2024melange} formulates heterogeneous GPU
allocation as cost-aware bin packing, achieving up to 77\% cost
reduction.  Our simulation tool~\cite{fleetsim2026} shows that the
cheapest GPU type depends on the workload in non-obvious ways:
A10G can beat H100 for concentrated-below workloads where
concurrency matters more than per-token speed.


% -----------------------------------------------------------------------------
\subsection{Energy and the 1/W Law}
\label{sec:pool:energy}
% -----------------------------------------------------------------------------

The 1/W law~\cite{onewlaw2026} has direct implications for pool
architecture: tokens per watt halves every time the context window
doubles, because of KV-cache concurrency limits.  At 64K
context, an H100 holds 16 sequences ($\tokW \approx 1.5$); at 4K,
256 sequences ($\tokW \approx 17.6$).

Routing topology is a stronger energy lever than GPU generation: two-pool routing gives $\sim$2.5$\times$ better
tok/W than homogeneous; upgrading H100$\to$B200 gives
$\sim$1.7$\times$.  The gains compose multiplicatively to
$\sim$4.25$\times$.  For MoE models like Qwen3-235B-A22B,
active-parameter weight streaming provides a third lever
($\sim$37.8\,tok/W at 8K context, 5.1$\times$ better than dense
Llama-3.1-70B).

Energy optimization therefore requires routing-pool co-design; a
homogeneous fleet cannot reach the energy frontier regardless of GPU
generation.


% -----------------------------------------------------------------------------
\subsection{Disaggregated Prefill/Decode and KV-Cache Topology}
\label{sec:pool:disagg}
% -----------------------------------------------------------------------------

Splitwise~\cite{patel2024splitwise} and
DistServe~\cite{zhong2024distserve} physically separate prefill and
decode, serving $4.48\times$ more requests at equivalent SLO.
Mooncake~\cite{qin2025mooncake} provides KV-cache-centric
disaggregation with hierarchical offloading (GPU HBM, CPU DRAM, SSD).
NVIDIA Dynamo~\cite{nvidia2026dynamo} enables dynamic scheduling with
KV-cache offloading across memory hierarchies.

Our inference-fleet-sim~\cite{fleetsim2026} models all three
topologies (monolithic, two-pool-routed, disaggregated), revealing
that disaggregation is not always beneficial: for concentrated-below
workloads (Archetype~1) at small $\boundary_{\text{short}}$, the
KV-cache transfer overhead exceeds the benefit.

vLLM's PagedAttention~\cite{kwon2023vllm} reduces KV-cache waste to
below 4\%, and prefix-cache-aware distributed
scheduling~\cite{llmd2026kvcache} delivers $57\times$ speedups.  But
for agentic workloads, standard LRU eviction is
pathological~\cite{vllm2026rfc37003}---context-aware retention
with session-affinity routing is needed.


